\section{Aproximaciones de diferencias finitas}

\begin{frame}
	\frametitle{\insertsection}

	El método numérico de diferencias finitas es universalmente
	aplicable para la solución de ecuaciones diferenciales.

	La idea básica es aproximar los cocientes diferenciales por un
	cociente de diferencias apropiado, por lo tanto, reducimos una
	ecuación diferencial a un sistema algebraico.

	Si $u\colon\mathbb{R}\to\mathbb{R}$ es 4-veces diferenciable,
	$x\in\mathbb{R}$ y $h>0$, entonces las siguientes aproximaciones
	de diferencias finitas son
	\begin{table}[ht!]
		\centering
		\begin{tabular}{ccc}
			\hline
			Derivada                       & Orden de precisión & Aproximación de diferencia finita                                                \\
			\hline
			                               & $1$                & $\dfrac{u_{i+1}-u_{i}}{h}$ o $\dfrac{u_{i}-u_{i-1}}{h}$                          \\[.2cm]
			$\displaystyle\diffp{u}{x}$    & $2$                & $\dfrac{u_{i+1}-u_{i-1}}{2h}$                                                    \\[.2cm]
			                               & $4$                & $\dfrac{u_{i-2}-8u_{i-1}+8u_{i+1}-u_{i+2}}{12h}$                                 \\[.2cm]
			\hline
			$\displaystyle\diffp[2]{u}{x}$ & $2$                & $\frac{u_{i-1}-2u_{i}+u_{i+1}}{h^2}$                                             \\[.2cm]
			                               & $4$                & $\frac{-u_{i-2} + 16u_{i-1} - 30u_{i} + 16u_{i+1} - u_{i+2}}{12h^2}$             \\[.2cm]
			\hline
			$\displaystyle\diffp[3]{u}{x}$ & $2$                & $\frac{-u_{i-2} + 2u_{i-1} - 2u_{i+1} + u_{i+2}}{2h^3}$                          \\[.2cm]
			                               & $4$                & $ \frac{u_{i-3} - 8u_{i-2} + 13u_{i-1} - 13u_{i+1} + 8u_{i+2} - u_{i+3}}{8h^3}$  \\[.2cm]
			\hline
			$\displaystyle\diffp[4]{u}{x}$ & $2$                & $\dfrac{u_{i-2}-4u_{i-1}+6u_{i}-4u_{i+1}+u_{i+2}}{h^4}$                          \\[.2cm]
			                               & $4$                & $\dfrac{-u_{i-3}+12u_{i-2}-39u_{i-1}+56u_{i}-39u_{i+1}+12u_{i+2}-u_{i+3}}{6h^4}$ \\[.2cm]
			\hline
		\end{tabular}
	\end{table}
\end{frame}

\begin{frame}
	\frametitle{\insertsection}

	Ahora usemos estas fórmulas de diferencias para formular algunos
	esquemas de diferencias para un ejemplo de problema de valores
	iniciales y de contorno para una ecuación del calor.
	\begin{exampleblock}{Ejemplo}
		Considere el siguiente problema de valor inicial / frontera
		\begin{equation}\label{eq:heatPDE}
			\begin{cases}\displaystyle
				\diffp{u}{t}=
				\nu\diffp[2]{u}{x}+
				f\left(x,t\right) &
				\text{en }\left(0,\pi\right)\times\left(0,T\right).   \\
				u=0               &
				\text{en }\left\{0,\pi\right\}\times\left(0,T\right). \\
				u=u_{0}           &
				\text{en }\left(0,\pi\right)\times\left\{t=0\right\}.
			\end{cases}
		\end{equation}
		Aquí $\nu>0$ es una constante dada, $f$ y $u_{0}$ son funciones
		continuas.
	\end{exampleblock}
	Para desarrollar un método de diferencias finitas, necesitamos
	introducir puntos de cuadrícula.
	Sean $N_{x}$ y $N_{t}$ números enteros positivos,
	$h_{x}=\dfrac{\pi}{N_{x}}$, $h_{t}=\dfrac{T}{Nt}$ y definir los
	puntos de partición
	\begin{align*}
		x_{j} & =
		jh_{x},\quad j=0,1,\dotsc,N_{x}. \\
		t_{m} & =
		mh_{t},\quad j=0,1,\dotsc,N_{t}.
	\end{align*}
	Un punto de la forma $\left(x_{j},t_{m}\right)$ se llama punto de
	cuadrícula y estamos interesados en calcular valores de solución
	aproximados en los puntos de la cuadrícula.
\end{frame}

\begin{frame}
	\frametitle{\insertsection}

	Usamos la notación $v^{m}_{j}$ para una aproximación a
	$u^{m}_{j}\equiv u\left(x_{j},t_{m}\right)$ calculado a partir de
	un esquema de diferencias finitas.
	Escribe $f^{m}_{j}=f\left(x_{j},t_{m}\right)$ y
	$r=\dfrac{\nu h_{t}}{h^{2}_{x}}$.
	Entonces podemos incorporar varios esquemas.
	El primer esquema es
	\begin{equation}\label{eq:heatPDEForward}
		\begin{cases}\displaystyle
			\dfrac{v^{m+1}_{j}-v^{m}_{j}}{h_{t}}=
			\nu
			\dfrac{v^{m}_{j+1}-2v^{m}_{j}+v^{m}_{j-1}}{h^{2}_{x}}+
			f^{m}_{j},                         &
			1\leq j\leq N_{x}-1, 0\leq m\leq N_{t}-1. \\
			v^{m}_{0}=v^{m}_{N_{x}}=0          &
			0\leq m\leq N_{t}.                        \\
			v^{0}_{j}=u_{0}\left(x_{j}\right), &
			0\leq j\leq N_{x}.
		\end{cases}
	\end{equation}
	Este esquema se obtiene discretizando la ecuación
	diferencial~\eqref{eq:heatPDE} en $x=x_{j}$ y $t=t_{m}$,
	reemplazando la derivada temporal con una diferencia directa y la
	segunda derivada espacial con una diferencia centrada de segundo
	orden.
	De ahí que se le llame esquema hacia adelante en tiempo y centrado
	en espacio.
	La ecuación de diferencias~\eqref{eq:heatPDEForward} se puede
	escribir como
	\begin{equation*}
		v^{m+1}_{j}=
		\left(1-2r\right)
		v^{m}_{j}+
		r\left(v^{m}_{j+1}+v^{m}_{j-1}\right)+
		h_{t}f^{m}_{j},\quad
		1\leq j\leq N_{x}-1,\quad 0\leq m\leq N_{t}-1.
	\end{equation*}
	Por lo tanto, una vez que se calcula la solución en el nivel de
	tiempo $t=t_{m}$, la solución en el siguiente nivel de tiempo
	$t=t_{m+1}$ se puede encontrar explícitamente.
	El esquema hacia adelante~\eqref{eq:heatPDEForward} es un método
	explícito.
\end{frame}

\begin{frame}
	\frametitle{\insertsection}

	Alternativamente, podemos reemplazar la derivada temporal con una
	diferencia hacia atrás y todavía use la diferencia centrada de
	segundo orden para la segunda derivada espacial.
	El esquema resultante es un esquema hacia atrás en el
	tiempo de espacio centrado:
	\begin{equation}\label{eq:heatPDEBackward}
		\begin{cases}\displaystyle
			\dfrac{v^{m}_{j}-v^{m-1}_{j}}{h_{t}}=
			\nu
			\dfrac{v^{m}_{j+1}-2v^{m}_{j}+v^{m}_{j-1}}{h^{2}_{x}}+
			f^{m}_{j},                         &
			1\leq j\leq N_{x}-1, 0\leq m\leq N_{t}-1. \\
			v^{m}_{0}=v^{m}_{N_{x}}=0,         &
			0\leq m\leq N_{t}.                        \\
			v^{0}_{j}=u_{0}\left(x_{j}\right), &
			0\leq j\leq N_{x}.
		\end{cases}
	\end{equation}
	La ecuación en diferencias~\eqref{eq:heatPDEBackward} se puede
	escribir como
	\begin{equation*}
		\left(1+2r\right)v^{m}_{j}-
		r\left(v^{m}_{j+1}+v^{m}_{j-1}\right)=
		v^{m-1}_{j}+h_{t}f^{m}_{j},\quad
		1\leq j\leq N_{x}-1,\quad
		1\leq m\leq N_{t},
	\end{equation*}
	que se complementa con la condición de frontera
	de~\eqref{eq:heatPDEBackward}.
	Por lo tanto, para encontrar la solución en el nivel de tiempo
	$t=t_{m}$ a partir de la solución en $t=t_{m-1}$,
	necesitamos resolver un sistema lineal tridiagonal de orden
	$N_{x}-1$.
	El esquema hacia atrás~\eqref{eq:heatPDEBackward} es un método
	implícito.

	En los dos métodos anteriores, aproximamos la ecuación diferencial
	en $x=x_{j}$ y $t=t_{m}$.
	También podemos considerar la ecuación diferencial en $x=x_{j}$ y
	$t=t_{m-\frac{1}{2}}$, aproximando la derivada del tiempo por una
	diferencia centrada:
	\begin{equation*}
		\diffp{u}{t}
		\left(x_{j},t_{m-\frac{1}{2}}\right)\approx
		\frac{u\left(x_{j},t_{m}\right)-u\left(x_{j},t_{m-1}\right)}{h_{t}}.
	\end{equation*}
\end{frame}

\begin{frame}
	\frametitle{\insertsection}

	Además, aproxima la segunda derivada espacial por la diferencia
	centrada de segundo orden:
	\begin{equation*}
		\diffp[2]{u}{x}
		\left(x_{j},t_{m-\frac{1}{2}}\right)\approx
		\frac{
			u\left(x_{j},t_{m}\right)-
			2u\left(x_{j},t_{m-\frac{1}{2}}\right)+
			u\left(x_{j},t_{m-1}\right)
		}{h^{2}_{x}},
	\end{equation*}
	y luego aproximar los valores de mitades de tiempo medio por
	promedios:
	\begin{equation*}
		u\left(x_{j},t_{m-\frac{1}{2}}\right)=
		\frac{1}{2}
		\left[
			u\left(x_{j},t_{m}\right)+
			u\left(x_{j},t_{m-1}\right)
			\right],
	\end{equation*}
	etc.
	Como resultado llegamos al esquema de Crank-Nicolson:
	\begin{equation}\label{eq:heatPDECrankNicolson}
		\begin{cases}\displaystyle
			\frac{v^{m}_{j}-v^{m-1}_{j}}{h_{t}}=
			\nu
			\frac{
				\left(v^{m}_{j+1}-2v^{m}_{j}+v^{m}_{j-1}\right)+
				\left(v^{m-1}_{j+1}-2v^{m-1}_{j}+v^{m-1}_{j-1}\right)
			}{2h^{2}_{x}}+
			f^{m-\frac{1}{2}}_{j},             &
			1\leq j\leq N_{x}-1, 1\leq m\leq N_{t}. \\
			v^{m}_{0}=v^{m}_{N_{x}}=0,         &
			0\leq m\leq N_{t}.                      \\
			v^{0}_{j}=u_{0}\left(x_{j}\right), &
			0\leq j\leq N_{x}.
		\end{cases}
	\end{equation}
	Aquí $f^{m-\frac{1}{2}}_{j}=f\left(x_{j},t_{m-\frac{1}{2}}\right)$,
	que puede ser reemplazado por $\left(f^{m}_{j}+f^{m-1}_{j}\right)/2$.
	La ecuación en diferencias~\eqref{eq:heatPDECrankNicolson} se puede
	reescribir como
	\begin{equation*}
		\left(1+r\right)v^{m}_{j}-
		\frac{r}{2}\left(v^{m}_{j+1}+v^{m}_{j-1}\right)=
		\left(1-r\right)v^{m-1}_{j}+
		\frac{r}{2}\left(v^{m-1}_{j+1}+v^{m-1}_{j-1}\right)+
		h_{t}f^{m-\frac{1}{2}}_{j},\quad
		1\leq j\leq N_{x}-1,\quad
		1\leq m\leq N_{t},
	\end{equation*}
	Vemos que el esquema de Crank-Nicolson también es un método
	implícito y en cada paso de tiempo necesitamos resolver un sistema
	lineal tridiagonal de orden $N_{x}-1$.
	Los tres esquemas derivados todos arriba parecen aproximaciones razonables
	para el problema de valor inicial frontera~\eqref{eq:heatPDE}.
\end{frame}

% Let us do some numerical experiments to see if these schemes indeed
% produce useful results.
% Let us use the forward scheme (6.1.8)–(6.1.10) and the backward scheme
% (6.1.12)–(6.1.14) to solve the problem (6.1.5)–(6.1.7) with ν = 1,
% f(x, t) = 0 and u0(x) = sin x.
% It can be verified that the exact solution is u(x, t) = e−t sin x.
% We consider numerical solution errors at t = 1.
% Results from the Crank-Nicolson scheme are qualitatively similar to
% those from the backward scheme but magnitudes of the error are
% smaller, and are thus omitted.
% We see that the Crank-Nicolson scheme is also an implicit method and at
% each time step we need to solve a tridiagonal linear system of order Nx−1.
% The three schemes derived above all seem reasonable approximations
% for the initial-boundary value problem (6.1.5)–(6.1.7).
% Let us do some numerical experiments to see if these schemes indeed
% produce useful results.
% Let us use the forward scheme (6.1.8)–(6.1.10) and the backward
% scheme (6.1.12)–(6.1.14) to solve the problem (6.1.5)–(6.1.7) with ν = 1,
% f(x, t) = 0 and u0(x) = sin x.
% It can be verified that the exact solution is u(x, t) = e−t sin x.
% We consider numerical solution errors at t = 1.
% Results from the Crank-Nicolson scheme are qualitatively similar to
% those from the backward scheme but magnitudes of the error are smaller,
% and are thus omitted.
